% !TeX program = xelatex
% ================================================================
% IEEEtran 中文会议论文模板（论文生成技能配套骨架）
% 编译：Tectonic -X compile main.tex（或 Overleaf XeLaTeX）
% 使用：全文替换占位内容；fig\ 子目录放图；保持每个 \label 唯一
% ================================================================

\documentclass[conference]{IEEEtran}
\IEEEoverridecommandlockouts

\usepackage{cite}
\usepackage{amsmath,amssymb,amsfonts}
\usepackage{graphicx}
\usepackage{textcomp}
\usepackage{xcolor}
\usepackage{booktabs}
\usepackage{multirow}
\usepackage{array}
\usepackage{url}
\usepackage{fontspec}
\usepackage{xeCJK}

\setCJKmainfont{Noto Serif CJK SC}
\setCJKsansfont{Noto Sans CJK SC}
\setCJKmonofont{Noto Sans Mono CJK SC}

\begin{document}

\title{论文标题：方法名与任务写全}

\author{
\IEEEauthorblockN{作者一, 作者二, 作者三}
\IEEEauthorblockA{
\textit{学院或实验室名称} \\
\textit{单位名称}\\
城市, 中国 \\
email@example.com}
}

\maketitle

\begin{abstract}
背景痛点两句。针对挑战一、挑战二、挑战三，本文提出方法名。方法模块按数据流各一句。在XX数据集上，本方法$mAP$达到XX.X\%，较最强基线提升X.X个百分点，较单模态/传统基线提升X.X个百分点。一句收束说明哪个机制被验证有效。
\end{abstract}

\begin{IEEEkeywords}
关键词一，关键词二，关键词三，关键词四，关键词五
\end{IEEEkeywords}

\section{引言}
% 应用领域 → 现有路线痛点 → 难点分析 → 数据集/基准背景 → 思想来源 → 贡献列表
本文主要贡献如下：
\begin{itemize}
    \item 贡献一：做了什么，解决什么。
    \item 贡献二：……
    \item 贡献三：……
\end{itemize}

\section{相关工作}
\subsection{方向一}
% 每小节末句收束到"还缺什么、本文补什么"
\subsection{方向二}

\section{数据集与问题定义}
\subsection{数据集}
\subsection{任务定义}
% 形式化：输入、输出、目标函数

\section{方法}
\subsection{总体结构}
% figure* 架构图
\subsection{模块一}
% 动机一句 → 公式 → 设计理由
\subsection{模块二}
\subsection{损失函数}

\section{实验与结果分析}
\subsection{实验设置}
\subsection{评价指标}
\subsection{总体性能比较}
% table* 主表：基线行用官方值，本方法行加粗；表注写明 ↑↓ 与加粗约定
\subsection{消融实验}
% 基线 < 单模块 < 组合 < 完整；参数量列自洽
\subsection{复杂环境/鲁棒性分析}
\subsection{定性分析}

\section{讨论}
\subsection{与最强基线的路线差异}
\subsection{局限性}

\section{结论}
% 方法两句复述 + 头条数字 + 未来工作

\begin{thebibliography}{00}
\bibitem{b1} 第一作者 et al., ``标题,'' \textit{期刊/会议}, 年.
\end{thebibliography}

\end{document}
